% This document was generated by an LLM.

\documentclass{essay}

\title{Oscillation: Measuring the Local Irregularity of Functions}
\author{}
\date{}

\begin{document}

\maketitle
\tableofcontents

%%%%%%%%%%%%%%%%%%%%%%%%%%%%%%%%%%%%%%%%%%%%%%%%%%%%%%%%%%%%%%%

\section{Introduction}

In the previous chapter we studied continuity and the various ways in which it
can fail. We learned that discontinuities come in several forms---removable,
jump, infinite, and oscillatory. While this classification is useful, it is
qualitative. It tells us \emph{what kind} of discontinuity is present but does
not measure \emph{how irregular} a function is near a point.

The concept of \emph{oscillation} provides exactly such a measurement.

Instead of asking whether a function is continuous, oscillation asks a more
refined question:

\begin{quote}
\emph{How much can the function vary inside an arbitrarily small neighborhood
of a point?}
\end{quote}

Remarkably, this simple idea connects several major topics in real analysis.

\begin{itemize}
\item Continuity can be characterized by zero oscillation.
\item Jump discontinuities correspond to positive oscillation.
\item Highly oscillatory functions exhibit maximal oscillation.
\item Riemann integrability can be described entirely in terms of oscillation.
\end{itemize}

Oscillation therefore serves as a bridge between continuity and integration.

%%%%%%%%%%%%%%%%%%%%%%%%%%%%%%%%%%%%%%%%%%%%%%%%%%%%%%%%%%%%%%%

\section{Variation on a Set}

Before discussing oscillation at a single point, it is useful to measure how
much a function varies on an arbitrary set.

\begin{definition}

Let $E\subseteq\mathbb{R}$ and suppose that $f$ is bounded on $E$.

The oscillation of $f$ on $E$ is

\[
\operatorname{osc}(f,E)
=
\sup_{x,y\in E}|f(x)-f(y)|.
\]

\end{definition}

This quantity measures the largest possible difference between function values
inside the set.

Equivalently,

\[
\operatorname{osc}(f,E)
=
\sup_{x\in E}f(x)
-
\inf_{x\in E}f(x).
\]

The second formula is often more convenient in analysis.

%%%%%%%%%%%%%%%%%%%%%%%%%%%%%%%%%%%%%%%%%%%%%%%%%%%%%%%%%%%%%%%

\section{Examples}

\begin{example}

If

\[
f(x)=5,
\]

then

\[
\operatorname{osc}(f,E)=0
\]

for every set $E$.

The function never changes value.

\end{example}

\begin{example}

Let

\[
f(x)=x
\]

on

\[
E=[1,4].
\]

Then

\[
\operatorname{osc}(f,E)=4-1=3.
\]

\end{example}

\begin{example}

For

\[
f(x)=\sin x
\]

on

\[
E=[0,2\pi],
\]

we obtain

\[
\operatorname{osc}(f,E)=2,
\]

since the function ranges from $-1$ to $1$.

\end{example}

%%%%%%%%%%%%%%%%%%%%%%%%%%%%%%%%%%%%%%%%%%%%%%%%%%%%%%%%%%%%%%%

\section{Local Oscillation}

The previous definition measures global variation over a set.

To understand continuity we must instead measure variation inside
progressively smaller neighborhoods.

For $\delta>0$, define

\[
I_\delta(x)
=
(x-\delta,x+\delta).
\]

The oscillation of $f$ inside this neighborhood is

\[
\omega_f(x,\delta)
=
\sup_{u,v\in I_\delta(x)}
|f(u)-f(v)|.
\]

As $\delta$ decreases, the neighborhood shrinks.

Global behavior gradually disappears, leaving only the local behavior near the
point.

%%%%%%%%%%%%%%%%%%%%%%%%%%%%%%%%%%%%%%%%%%%%%%%%%%%%%%%%%%%%%%%

\section{Oscillation at a Point}

We are now ready for the central definition.

\begin{definition}

Suppose $f$ is bounded in some neighborhood of $x$.

The oscillation of $f$ at $x$ is

\[
\omega_f(x)
=
\lim_{\delta\to0}
\omega_f(x,\delta).
\]

\end{definition}

Why does this limit exist?

Observe that if

\[
0<\delta_1<\delta_2,
\]

then

\[
I_{\delta_1}(x)
\subseteq
I_{\delta_2}(x).
\]

Therefore,

\[
\omega_f(x,\delta_1)
\le
\omega_f(x,\delta_2).
\]

Thus the oscillation cannot increase as the neighborhood becomes smaller.

The function

\[
\delta
\longmapsto
\omega_f(x,\delta)
\]

is monotone decreasing.

Since it is also bounded below by zero, the limit always exists (possibly equal
to zero).

%%%%%%%%%%%%%%%%%%%%%%%%%%%%%%%%%%%%%%%%%%%%%%%%%%%%%%%%%%%%%%%

\section{Interpretation}

Oscillation provides a numerical measure of local irregularity.

If

\[
\omega_f(x)=0,
\]

then every sufficiently small neighborhood contains function values that are
arbitrarily close together.

The function becomes nearly constant when viewed under sufficient magnification.

If

\[
\omega_f(x)
\]

is positive, then no matter how closely we zoom in, there remain points whose
function values differ by at least some fixed amount.

Thus oscillation measures how stubbornly the function refuses to settle down.

%%%%%%%%%%%%%%%%%%%%%%%%%%%%%%%%%%%%%%%%%%%%%%%%%%%%%%%%%%%%%%%

\section{Examples}

\subsection{Continuous Functions}

Every continuous function satisfies

\[
\omega_f(x)=0
\]

at every point.

Indeed, continuity states that nearby function values become arbitrarily close
to the function value at the center point.

%%%%%%%%%%%%%%%%%%%%%%%%%%%%%%%%%%%%%%%%%%%%%%%%%%%%%%%%%%%%%%%

\subsection{Jump Discontinuities}

Consider

\[
f(x)=
\begin{cases}
0,&x<0,\\
1,&x\ge0.
\end{cases}
\]

Every neighborhood of the origin contains points whose values are $0$ and $1$.

Hence

\[
\omega_f(0)=1.
\]

Away from the origin,

\[
\omega_f(x)=0.
\]

%%%%%%%%%%%%%%%%%%%%%%%%%%%%%%%%%%%%%%%%%%%%%%%%%%%%%%%%%%%%%%%

\subsection{The Floor Function}

The floor function

\[
\lfloor x\rfloor
\]

has oscillation

\[
1
\]

at every integer.

Elsewhere,

\[
\omega_f(x)=0.
\]

%%%%%%%%%%%%%%%%%%%%%%%%%%%%%%%%%%%%%%%%%%%%%%%%%%%%%%%%%%%%%%%

\subsection{A Removable Discontinuity}

Consider

\[
f(x)=
\begin{cases}
1,&x=0,\\
0,&x\neq0.
\end{cases}
\]

Every neighborhood of the origin contains both values $0$ and $1$.

Consequently,

\[
\omega_f(0)=1.
\]

Although the discontinuity is removable, its oscillation is still positive.

%%%%%%%%%%%%%%%%%%%%%%%%%%%%%%%%%%%%%%%%%%%%%%%%%%%%%%%%%%%%%%%

\subsection{Rapid Oscillations}

Now consider

\[
f(x)
=
\sin\left(\frac1x\right),
\qquad x\neq0.
\]

Every neighborhood of the origin contains points where the function equals
$1$ and points where it equals $-1$.

Therefore

\[
\omega_f(0)=2.
\]

This is the largest possible oscillation for a function whose values lie between
$-1$ and $1$.

%%%%%%%%%%%%%%%%%%%%%%%%%%%%%%%%%%%%%%%%%%%%%%%%%%%%%%%%%%%%%%%

\subsection{Dirichlet's Function}

Define

\[
f(x)=
\begin{cases}
1,&x\in\mathbb{Q},\\
0,&x\notin\mathbb{Q}.
\end{cases}
\]

Every interval contains both rational and irrational numbers.

Hence

\[
\omega_f(x)=1
\]

for every real number.

The function is therefore discontinuous everywhere.

%%%%%%%%%%%%%%%%%%%%%%%%%%%%%%%%%%%%%%%%%%%%%%%%%%%%%%%%%%%%%%%

\subsection{Thomae's Function}

Thomae's function is defined by

\[
f(x)=
\begin{cases}
1/q,
&
x=p/q
\text{ in lowest terms},\\
0,
&
x\notin\mathbb{Q}.
\end{cases}
\]

This remarkable example satisfies

\[
\omega_f(x)=0
\]

at every irrational point,

but

\[
\omega_f(x)>0
\]

at every rational point.

It is continuous exactly at the irrational numbers.

%%%%%%%%%%%%%%%%%%%%%%%%%%%%%%%%%%%%%%%%%%%%%%%%%%%%%%%%%%%%%%%

\section{Characterizing Continuity}

Oscillation provides one of the cleanest characterizations of continuity.

\begin{theorem}

A function is continuous at a point $x$ if and only if

\[
\omega_f(x)=0.
\]

\end{theorem}

Thus continuity is equivalent to saying that the function becomes arbitrarily
flat when viewed under sufficiently high magnification.

%%%%%%%%%%%%%%%%%%%%%%%%%%%%%%%%%%%%%%%%%%%%%%%%%%%%%%%%%%%%%%%

\section{Oscillation and Upper and Lower Sums}

Oscillation plays a central role in Riemann integration.

Consider a subinterval

\[
I=[a,b].
\]

The oscillation of $f$ on $I$ is

\[
\sup_I f-\inf_I f.
\]

If the oscillation is small, then every value of the function on that interval
is nearly the same.

Consequently,

\begin{itemize}
\item the upper rectangle,
\item the lower rectangle,
\end{itemize}

have nearly equal heights.

Their contribution to the difference between the upper and lower Riemann sums is
therefore small.

Large oscillation produces large discrepancies between upper and lower sums.

%%%%%%%%%%%%%%%%%%%%%%%%%%%%%%%%%%%%%%%%%%%%%%%%%%%%%%%%%%%%%%%

\section{Oscillation and Riemann Integrability}

The importance of oscillation culminates in one of the central theorems of real
analysis.

\begin{theorem}[Lebesgue's Criterion for Riemann Integrability]

A bounded function on a closed interval is Riemann integrable if and only if the
set of points where

\[
\omega_f(x)>0
\]

has measure zero.

\end{theorem}

Although measure theory lies beyond the scope of this chapter, the theorem has a
simple interpretation.

A bounded function is Riemann integrable provided its discontinuities occupy a
negligibly small portion of the interval.

This explains why piecewise continuous functions are always Riemann integrable.

%%%%%%%%%%%%%%%%%%%%%%%%%%%%%%%%%%%%%%%%%%%%%%%%%%%%%%%%%%%%%%%

\section{Oscillation Versus Other Measures of Regularity}

Oscillation measures only local variation.

Other notions of regularity describe different aspects of a function.

\begin{center}
\begin{tabular}{ll}
Continuity & Local predictability\\
Uniform continuity & Global control of continuity\\
Differentiability & Local linear approximation\\
Lipschitz continuity & Bounded rate of change\\
Bounded variation & Total accumulated variation\\
Oscillation & Local spread of function values
\end{tabular}
\end{center}

These concepts complement one another rather than compete.

%%%%%%%%%%%%%%%%%%%%%%%%%%%%%%%%%%%%%%%%%%%%%%%%%%%%%%%%%%%%%%%

\section{A Geometric View}

Imagine repeatedly zooming into the graph of a function.

For a continuous function, every magnification reveals a graph that becomes
flatter and flatter.

The oscillation tends toward zero.

For a jump discontinuity, every zoom still reveals two separated branches.

The oscillation never disappears.

For

\[
\sin\left(\frac1x\right),
\]

every magnification uncovers new oscillations.

No matter how closely we inspect the graph, it never resembles a continuous
curve.

Oscillation therefore captures exactly what our geometric intuition observes.

%%%%%%%%%%%%%%%%%%%%%%%%%%%%%%%%%%%%%%%%%%%%%%%%%%%%%%%%%%%%%%%

\section{Summary}

Oscillation transforms continuity from a binary concept into a quantitative one.

Instead of merely asking whether a function is continuous, it measures how much
the function continues to vary inside arbitrarily small neighborhoods.

This single idea unifies several seemingly unrelated topics.

It explains why continuous functions have stable local behavior, why jump
discontinuities remain visible under magnification, why rapidly oscillating
functions resist limiting values, and why Riemann integration depends on the
distribution of discontinuities rather than their mere existence.

For these reasons, oscillation occupies a central position in real analysis,
providing one of the most elegant connections between continuity, integration,
and the local geometry of functions.

\end{document}
